\documentclass{standalone}
\usepackage{circuitikz}
\usetikzlibrary{calc}
\usetikzlibrary{arrows.meta}
\definecolor{circuitnoteblue}{HTML}{4C8EDA}
\definecolor{circuitnotered}{HTML}{E5534B}
\definecolor{circuitnotemark}{HTML}{FE00FE}
\begin{document}
\begin{circuitikz}[american, line width=0.8pt]
\ctikzset{bipoles/length=1.2cm}
\ctikzset{grounds/scale=1.36}
\foreach \x in {0,2,4,6,8,10} {\foreach \y in {0,-2,-4} {\fill[gray, opacity=0.35] (\x,\y) circle (1.1pt);}}
\node[gray, font=\scriptsize] at (0,0.84) {1};
\node[gray, font=\scriptsize] at (2,0.84) {2};
\node[gray, font=\scriptsize] at (4,0.84) {3};
\node[gray, font=\scriptsize] at (6,0.84) {4};
\node[gray, font=\scriptsize] at (8,0.84) {5};
\node[gray, font=\scriptsize] at (10,0.84) {6};
\node[gray, font=\scriptsize, anchor=east] at (-0.84,0) {a};
\node[gray, font=\scriptsize, anchor=east] at (-0.84,-2) {b};
\node[gray, font=\scriptsize, anchor=east] at (-0.84,-4) {c};
\coordinate (a1) at (0,0);
\coordinate (a3) at (4,0);
\coordinate (c3) at (4,-4);
\coordinate (a4) at (6,0);
\draw (a1) node[ocirc]{} node[above left]{IN}; % line 3
\draw (a1) to[R, l_=$R_{1}$, a^=$10\,\mathrm{k}\Omega$] (a3); % line 4
\draw (a3) to[C, l_=$C_{1}$, a^=$100\,\mathrm{n}\mathrm{F}$] (c3); % line 5
\draw (a4) node[ocirc]{} node[above left]{OUT}; % line 6
\node[ground] at (c3) {}; % line 7
\draw (a3) -- (a4); % line 9
\node[circ] at (a3) {};
\draw[circuitnoteblue, dashed, rounded corners=3pt] (-0.7,-4.7) rectangle (4.7,0.7); % line 11
\path (0.222,-6.593) rectangle (3.778,-5.407); % line 12
\node[anchor=west, inner sep=0, circuitnotemark, font=\large] at (2,-6) {X}; % line 12
\draw[circuitnotered, -{Stealth[length=2.2mm]}] (10,-4) -- (2.805,-0.402); % line 13
\path (10,-4.395) rectangle (11.863,-3.605); % line 14
\node[anchor=west, inner sep=0, circuitnotemark, font=\footnotesize] at (10,-4) {X}; % line 14
\path (12,-9.842) rectangle (23.811,0.593); % line 15
\node[anchor=west, inner sep=0, circuitnotemark, font=\large] at (12,0) {X}; % line 15
\node[anchor=west, inner sep=0, circuitnotemark, font=\large] at (12,-0.487) {X}; % line 15
\node[anchor=west, inner sep=0, circuitnotemark, font=\large] at (12,-0.974) {X}; % line 15
\node[anchor=west, inner sep=0, circuitnotemark, font=\large] at (12,-1.46) {X}; % line 15
\node[anchor=west, inner sep=0, circuitnotemark, font=\large] at (12,-1.947) {X}; % line 15
\node[anchor=west, inner sep=0, circuitnotemark, font=\large] at (12,-2.434) {X}; % line 15
\node[anchor=west, inner sep=0, circuitnotemark, font=\large] at (12,-2.921) {X}; % line 15
\node[anchor=west, inner sep=0, circuitnotemark, font=\large] at (12,-3.408) {X}; % line 15
\node[anchor=west, inner sep=0, circuitnotemark, font=\large] at (12,-3.895) {X}; % line 15
\node[anchor=west, inner sep=0, circuitnotemark, font=\large] at (12,-4.381) {X}; % line 15
\node[anchor=west, inner sep=0, circuitnotemark, font=\large] at (12,-4.868) {X}; % line 15
\node[anchor=west, inner sep=0, circuitnotemark, font=\large] at (12,-5.355) {X}; % line 15
\node[anchor=west, inner sep=0, circuitnotemark, font=\large] at (12,-5.842) {X}; % line 15
\node[anchor=west, inner sep=0, circuitnotemark, font=\large] at (12,-6.329) {X}; % line 15
\node[anchor=west, inner sep=0, circuitnotemark, font=\large] at (12,-6.816) {X}; % line 15
\node[anchor=west, inner sep=0, circuitnotemark, font=\large] at (12,-7.302) {X}; % line 15
\node[anchor=west, inner sep=0, circuitnotemark, font=\large] at (12,-7.789) {X}; % line 15
\node[anchor=west, inner sep=0, circuitnotemark, font=\large] at (12,-8.276) {X}; % line 15
\node[anchor=west, inner sep=0, circuitnotemark, font=\large] at (12,-8.763) {X}; % line 15
\node[anchor=west, inner sep=0, circuitnotemark, font=\large] at (12,-9.25) {X}; % line 15
\node[anchor=south west, inner sep=2pt, circuitnotemark, font=\large\bfseries] (circuittitle) at (current bounding box.north west) {X};
\path (circuittitle.south west) rectangle ++(3.874,0.593);
\end{circuitikz}
\end{document}
